% FedAuditor project walkthrough for the guide, in the app's "deep ink" design.
% Build in this folder:  pdflatex FedAuditor_Walkthrough ; pdflatex FedAuditor_Walkthrough
% Console outputs: `python -m fedauditor.walkthrough` (results/fedauditor/walkthrough/)
% Tables: `python -m fedauditor.make_figures` (results/fedauditor/figures/)
% Screenshots + ink art: docs/design/ (make_ink_art.py; screens/ captured from the running app)
\documentclass[11pt,a4paper]{article}
\usepackage[margin=2.2cm,top=2.4cm,bottom=2.4cm]{geometry}
\usepackage[T1]{fontenc}
\usepackage[scaled=0.95]{helvet}
\renewcommand{\familydefault}{\sfdefault}
\usepackage{amsmath}
\usepackage{booktabs}
\usepackage[table]{xcolor}
\usepackage{graphicx}
\usepackage{textcomp}
\usepackage{fancyvrb}
\usepackage{float}
\usepackage{placeins}
\usepackage{enumitem}
\usepackage[most]{tcolorbox}
\usepackage{tikz}
\usepackage{titlesec}
\usepackage{caption}
\usepackage{fancyhdr}
\usepackage[hidelinks]{hyperref}

% ---- Deep-ink palette (same tokens as static/css/style.css) ----
\definecolor{inkbase}{HTML}{04050C}
\definecolor{fog}{HTML}{EEF0F6}
\definecolor{muted}{HTML}{B9BECF}
\definecolor{inkcyan}{HTML}{67E8F9}
\definecolor{inkviolet}{HTML}{A78BFA}
\definecolor{inkmagenta}{HTML}{F0ABFC}
\definecolor{glass}{HTML}{0E1220}
\definecolor{glasstitle}{HTML}{161B2C}
\definecolor{glassline}{HTML}{2A3045}
\definecolor{mint}{HTML}{6EE7B7}
\definecolor{good}{HTML}{4ADE80}

\pagecolor{inkbase}
\color{fog}
\arrayrulecolor{glassline}

\setlist{itemsep=3pt,topsep=4pt,label={\color{inkcyan}\textbullet}}
\setlength{\parskip}{5pt}
\setlength{\parindent}{0pt}
% Let tables share pages with text instead of getting float-only pages
\renewcommand{\topfraction}{0.9}
\renewcommand{\bottomfraction}{0.8}
\renewcommand{\textfraction}{0.07}
\renewcommand{\floatpagefraction}{0.85}
\setcounter{topnumber}{4}
\setcounter{totalnumber}{6}
\renewcommand{\baselinestretch}{1.08}

% Headings: medium weight, fog-white, cyan section numbers
\titleformat{\section}{\color{fog}\Large\bfseries}{\color{inkcyan}\thesection}{0.8em}{}
\titlespacing*{\section}{0pt}{18pt}{8pt}
\captionsetup{labelfont={color=inkcyan,bf},textfont={color=muted},font=small}

% Page frame: muted running header and page number
\pagestyle{fancy}
\fancyhf{}
\renewcommand{\headrulewidth}{0pt}
\fancyhead[L]{\color{muted}\footnotesize FedAuditor \textperiodcentered{} Project walkthrough}
\fancyhead[R]{\color{muted}\footnotesize\thepage}

\newcommand{\res}{../../results/fedauditor}
\newcommand{\design}{../design}

% Console output: the app's deepest surface, mint text
\newcommand{\out}[1]{\begin{tcolorbox}[enhanced,colback=glass,colframe=glassline,coltext=fog,boxrule=0.5pt,arc=6pt,
  left=6pt,right=6pt,top=3pt,bottom=3pt,
  title={\footnotesize Console output: \texttt{#1}},coltitle=muted,colbacktitle=glasstitle,fonttitle=\footnotesize]
  \VerbatimInput[fontsize=\scriptsize,formatcom=\color{mint}]{\res/walkthrough/#1}\end{tcolorbox}}

% Screenshot in a glass frame
\newcommand{\shot}[2]{\begin{figure}[H]\centering
  \begin{tcolorbox}[enhanced,colback=glass,colframe=glassline,boxrule=0.5pt,arc=8pt,
    left=3pt,right=3pt,top=3pt,bottom=3pt,width=0.92\textwidth]
  \includegraphics[width=\linewidth]{\design/screens/#1}\end{tcolorbox}
  \caption{#2}\end{figure}}

% Glass card for summaries
\newtcolorbox{glasscard}[1]{enhanced,colback=glass,colframe=glassline,coltext=fog,boxrule=0.5pt,arc=10pt,
  left=10pt,right=10pt,top=8pt,bottom=8pt,title={#1},fonttitle=\bfseries,
  coltitle=inkbase,colbacktitle=fog,attach boxed title to top left={xshift=10pt,yshift=-9pt},
  boxed title style={arc=7pt,boxrule=0pt,colback=fog}}

\hypersetup{colorlinks=true,linkcolor=fog,urlcolor=inkcyan}

\begin{document}

% ======================= Cover =======================
\thispagestyle{empty}
\begin{tikzpicture}[remember picture,overlay]
  \node[inner sep=0pt] at (current page.center)
    {\includegraphics[height=\paperheight]{\design/ink_cover.png}};
  \fill[inkbase,opacity=0.18] (current page.south west) rectangle (current page.north east);
  \node[anchor=north west,text=fog,font=\large\bfseries] at ([xshift=2.2cm,yshift=-2.4cm]current page.north west)
    {FedHealth \textcolor{muted}{\textperiodcentered{} FedAuditor}};
  \node[anchor=west,text width=15cm,align=left] at ([xshift=2.2cm,yshift=3.2cm]current page.west) {%
    \tcbox[enhanced,colback=glass,colframe=glassline,boxrule=0.5pt,arc=9pt,
      left=8pt,right=8pt,top=3pt,bottom=3pt,on line]{\color{muted}\small Final-year project \textperiodcentered{} Guide review}\\[14pt]
    {\color{fog}\fontsize{38}{44}\selectfont\bfseries\spaceskip=0.28em plus 0.05em Audit Federated\\ Models Without\\ Seeing Patient Data\par}
    \vspace{12pt}
    {\color{muted}\large From FedHealth to a privacy-budgeted, fact-checked LLM auditor\\ for federated clinical models. Every step and its output.\par}};
  \node[anchor=south west,text=muted,font=\small] at ([xshift=2.2cm,yshift=2.2cm]current page.south west)
    {November 2026 \quad\textperiodcentered\quad Records stay local; only noisy aggregates leave a site.};
\end{tikzpicture}
\newpage

\begin{glasscard}{Summary}
\textbf{\color{inkcyan}Problem.} In federated learning the coordinator never sees patient data, so it cannot tell
when one hospital or region has a fault (wrong labels, broken exports, an under-diagnosed subgroup,
a malicious update). Pooled accuracy does not show these faults.

\textbf{\color{inkcyan}What we built.} An LLM agent that audits a federated model using only noisy, privacy-budgeted
aggregate tools. It never sees a patient record, and every number in its report is checked against
the tool log by a fact-checker. To test it we built \textbf{FedFault}, a benchmark with 8 planted
fault types $\times$ 3 severities $\times$ 5 seeds, on real CDC BRFSS data with 52 US states as clients.

\textbf{\color{inkcyan}Result.} Over 125 test scenarios, Claude Haiku~5.5 finds the faulty clients with precision 0.97
and recall 0.89, names the fault in 91\% of scenarios, and 2{,}262 of its 2{,}264 claims are verified;
the whole evaluation cost US\$1.16. A tuned rule baseline scores 0.99 / 0.88 / 87\%. The agent is
much better at subgroup bias. A 6-page IEEE conference paper and a redesigned web app are ready.
\end{glasscard}

\vspace{6pt}
{\color{fog}\tableofcontents}
\newpage

% =================================================================================================
\section{Starting point: FedHealth}

The original project (FedHealth) trained a logistic regression with FedAvg over \textbf{3 simulated
age-cohort clients} on the PIMA dataset (768 records), with a Flask dashboard and a measurement layer
(calibration, subgroup fairness, explanation stability, client Shapley values). The guide's feedback
was that it was ``very basic''. We kept all of it and built on top of it: the measurement modules
became the agent's tools.

\shot{0_original_dashboard.jpg}{The FedHealth dashboard (PIMA, 768 records, 3 age-cohort clients), shown in the new design.}

\textbf{What changed}, in five phases:
\begin{center}\small
\begin{tabular}{lll}
\toprule
Phase & Work & Main output \\
\midrule
1 & Real data and federation & BRFSS 2023, 52 state clients, LR/MLP/XGBoost \\
2 & FedFault benchmark & 8 faults $\times$ 3 severities $\times$ 5 seeds + clean runs, ground truth \\
3 & Agent, tools, fact-checker & 8 noisy budgeted tools, Claude agent, claim checker \\
4 & Experiments & 125 audits per auditor, rule baseline, prompt ablation, open models \\
5 & Paper and dashboard & 6-page IEEE paper, FedAuditor page in the Flask app \\
\bottomrule
\end{tabular}
\end{center}

All outputs in Sections 2--5 below were produced by running the project's own scripts; the
console boxes are copied unedited from \texttt{python -m fedauditor.walkthrough}, which re-runs the
whole pipeline live (free, no API calls).

% =================================================================================================
\section{Phase 1: real data and a real federation}

We replaced PIMA (768 rows, 3 artificial clients) with the CDC Behavioral Risk Factor Surveillance
System 2023: about 433{,}000 survey respondents, recoded to a diabetes outcome with 17 risk factors.
Each US state or territory is one client, so the federation is naturally non-identical and very
unbalanced. Missing-value imputation and feature scaling are computed in a federated way.

\out{step1.txt}

The federation runs with the Flower framework for 20 rounds of FedAvg. Three model families are
supported: logistic regression (LR), a small neural network (MLP) and XGBoost (tree bagging).

\out{step2.txt}

\begin{center}\small
\begin{tabular}{lccc}
\toprule
Model & Federated AUC & Central AUC (pooled data) & Per-state AUC \\
\midrule
Logistic regression & 0.821 & 0.821 & 0.76--0.87 \\
XGBoost & 0.826 & 0.826 & \\
MLP & 0.816 & 0.828 (still rising) & \\
\bottomrule
\end{tabular}
\end{center}
Federated training matches central training, so the federation itself is sound before any fault is
planted.

% =================================================================================================
\section{Phase 2: the FedFault benchmark}

To measure an auditor we need faults with a known answer. Each scenario plants one fault in 3 clients
(about 5\% of the federation), chosen among the larger clients and kept the same at every severity.

\begin{center}\small
\begin{tabular}{p{3.6cm}p{8.6cm}c}
\toprule
Fault & What is planted & Low / Med / High \\
\midrule
Label noise & fraction of labels flipped & 0.1 / 0.2 / 0.3 \\
Drift & from mid-training, top-3 risk features shift by $k$ SD & 0.5 / 1 / 2 \\
Poisoning & update replaced by $-k$ times itself & 1 / 3 / 10 \\
Free-rider & fraction of rounds the client sends back the global model & 0.3 / 0.6 / 1.0 \\
Prevalence shift & only $1/k$ of negatives kept & 2 / 4 / 8 \\
Subgroup bias & fraction of positives in one subgroup recorded as negative & 0.3 / 0.6 / 0.9 \\
Missing-data artefact & fraction of rows with two features exported as a sentinel code & 0.1 / 0.25 / 0.5 \\
Explanation instability & noise of $k\times$ the update norm added to every upload & 1 / 3 / 10 \\
\bottomrule
\end{tabular}
\end{center}

Seed 10 is used only for development; seeds 20--24 are the test set: 120 faulty + 5 clean = 125
scenarios per model, run for LR, MLP and XGBoost (375 federations).

\out{step3.txt}

The key observation, over all 375 runs: \textbf{pooled AUC hides almost every fault} (Table~\ref{tab:pooled}).
Only high-severity poisoning is visible. Everything else changes AUC by at most 0.001, less than the
0.008 that AUC varies between clean runs with different seeds. This is why an auditor is needed.

\input{\res/figures/pooled_auc.tex}

% =================================================================================================
\section{Phase 3: privacy-budgeted tools, the agent and the fact-checker}

\textbf{Tools.} The auditor can never read records. It calls 8 tools that compute statistics at each
client and release only noisy aggregates:
\begin{itemize}
\item \texttt{federation\_overview} (free): clients, sizes, features, sensitive attributes;
\item \texttt{global\_performance}, \texttt{client\_performance}: AUC, calibration, prevalence per client;
\item \texttt{client\_updates}: how large each client's update is and whether it points the same way as the others;
\item \texttt{feature\_summary}, \texttt{subgroup\_outcomes}, \texttt{drift\_by\_round}, \texttt{attribution\_churn}.
\end{itemize}
Each paid call adds Laplace noise of scale $c\,m/(n\varepsilon)$ and spends $\varepsilon$ from a budget of
20 per audit. Because noise grows as clients get smaller, every table carries \emph{noise-adjusted}
flags, so that small states are not flagged just for being noisy.

\out{step4.txt}

\textbf{Agent.} A Claude model runs a tool-use loop: it screens all clients, drills into suspicious ones,
and submits a JSON report. Each finding names a client, fault type, onset round and confidence, and
lists \emph{claims}: (tool call id, metric, value, client/round/feature/group).

\textbf{Fact-checker.} Every claim is looked up in the tool log; it is accepted only if the cited call
has a row with the same keys and a value within $0.005 + 2\%$. \textbf{Rule baseline.} A non-LLM
auditor with the same tools, budget and report format (fixed plan + thresholds), so we can see what
the LLM adds. Below, the rule auditor audits the poisoning scenario from Step 3, end to end:

\out{step5.txt}

The first live LLM audit (Claude Opus~5.5, same scenario) also flagged MD, MI and MT as poisoning
from round 1: precision, recall and type accuracy 1.0, 7/7 claims verified, 15 tool calls, about
US\$0.31.

% =================================================================================================
\section{Phase 4: experiments}

We audited all 125 LR test scenarios with the rule baseline and with Claude Haiku~5.5 (two prompt
versions), scoring every report against the ground truth.

\out{step6.txt}

\input{\res/figures/main_table.tex}

\input{\res/figures/recall_heatmap_dark.tex}
\input{\res/figures/type_confusion_dark.tex}
\input{\res/figures/behaviour.tex}
\FloatBarrier

\textbf{Findings.}
\begin{itemize}
\item The agent and the rules find faulty clients equally well (recall 0.89 vs 0.88); the agent
  names the fault better (0.91 vs 0.87) and dates it better (0.24 vs 0.35 rounds).
\item \textbf{Subgroup bias}: the agent finds every affected client at medium severity; the rules find 60\%
  (Table~\ref{tab:recall}).
\item The agent is weaker on the faintest drift and free-riding. Nobody detects the lowest severity of
  subgroup bias or explanation instability: the signal is below the privacy noise.
\item \textbf{Prompt v5} added one rule (``check local AUC before calling poisoning''): label noise
  mis-typed as poisoning dropped from 12/15 scenarios to 0/15, and type accuracy rose 0.79 $\to$ 0.91
  (Table~\ref{tab:confusion}).
\item Both agent false alarms are Guam, the smallest client, at confidence exactly 0.5.
\item The fact-checker rejected 9 claims (prompt v4) and 2 (v5); two v4 claims were genuinely wrong numbers
  (O/E 7.38 claimed vs 1.81 logged). Cost: under one cent per audit.
\end{itemize}


\input{\res/figures/rules_models.tex}

\textbf{Open-weight models} (Qwen3 8B and 30B, run free on Kaggle/locally with the same tools and prompt)
failed for model reasons: the 8B model fixed on the smallest state and never screened; the 30B
thinking model took 16--29 minutes per audit and read noise as faults; the 30B instruct model caught
poisoning but over-flagged weak signals and mis-cited numbers, which the fact-checker caught. This is
reported in the paper as a finding.

% =================================================================================================
\section{Phase 5: dashboard and paper}

\textbf{Dashboard.} A new \emph{FedAuditor} page in the Flask app shows any saved audit. Every claim is a
link; clicking it opens the tool call it cites and highlights the exact logged row.

\shot{1_report_top.jpg}{FedAuditor page: planted fault, flagged clients, type, claims verified (25/25), privacy budget, and the agent's summary.}
\shot{2_claims.jpg}{One finding with its claims. Each claim shows the claimed value, the logged value, the cited call and the check result.}
\shot{3_cited_row.jpg}{After clicking claim \texttt{c6}: the cited \texttt{subgroup\_outcomes} call opens and the row for MA is highlighted (O/E ratio 0.0, $z_\text{noise}=-8.98$).}
\shot{4_figures.jpg}{Benchmark results shown on the same page: recall heatmap and fault-type confusion.}

\textbf{Paper.} A 6-page IEEE conference paper (\texttt{paper/fedauditor.pdf}): introduction, related work
(29 references, including the four closest papers FedAgentBench, Agentic-FL, FL-MAESTRO and
ModelAuditor, none of which audits a federated model), problem setting, FedFault, method, experiments,
discussion, limitations. All its tables are generated by a script and checked against the saved results.

\textbf{New design.} The whole web app now shares one look, taken from a fluid ``ink in water'' hero:
a near-black base, fog-white type, translucent glass cards, white pill buttons and a
cyan\,$\to$\,violet\,$\to$\,magenta accent band. The landing and sign-in pages run a live WebGL ink
simulation; data pages stay calm so tables and charts remain easy to read. Layouts and features are unchanged.

\shot{hero_landing.jpg}{Landing page: live fluid-ink background, word-by-word heading, glass call-to-action bar.}
\shot{login.jpg}{Sign-in: the same ink behind a glass card.}
\shot{predict.jpg}{Diagnostic form with a prediction: glass inputs, white pill toggles, contribution bars.}

% =================================================================================================
\section{How to reproduce}

All commands run from \texttt{Final-Year/} with the project's virtual environment:
\begin{center}\footnotesize
\begin{tabular}{ll}
\toprule
Step & Command \\
\midrule
Download and recode BRFSS & \texttt{python -m fedauditor.prepare\_brfss -{}-year 2023} \\
Clean baseline & \texttt{python -m fedauditor.run\_baseline -{}-dataset brfss} \\
FedFault grid (375 runs) & \texttt{python -m fedauditor.run\_fedfault -{}-models lr mlp xgb} \\
One LLM audit (paid) & \texttt{python -m fedauditor.run\_audit -{}-fault poisoning -{}-severity low} \\
Audit grid & \texttt{python -m fedauditor.run\_audit\_grid -{}-agents rules claude} \\
Compare auditors & \texttt{python -m fedauditor.compare\_agents -{}-model lr} \\
Tables and figures & \texttt{python -m fedauditor.make\_figures} \\
This walkthrough (free) & \texttt{python -m fedauditor.walkthrough} \\
Web app & \texttt{python app.py} $\to$ \url{http://127.0.0.1:5000/fedauditor} \\
Tests & \texttt{python -m unittest discover -s tests} (78 tests, all pass) \\
\bottomrule
\end{tabular}
\end{center}

\section{Limitations and next steps}
\begin{itemize}
\item The privacy budget is a calibrated accounting device, not a formal differential-privacy proof
  for every statistic (exact for means and rates, heuristic for AUC and update statistics).
\item LLM audits were run on LR only (API budget); the rule baseline covers LR, MLP and XGBoost.
\item One dataset; five clean runs, so false-alarm rates are coarse.
\item The fact-checker checks numbers, not wording: in one report the agent called the \texttt{Sex}$=1$
  group ``women'' although 1 means men.
\item Next: formal privacy mechanism, LLM audits on MLP/XGBoost, a second dataset (eICU) for the
  journal version, and checking qualitative claims.
\end{itemize}

\end{document}
